\scanpage{17}
\mspage{31}

En particulier, on obtient:

\uline{Corollaire} \inserted{4.}12. Soient $A\longrightarrow B$ un homomorphisme local d'anneaux locaux
noethériens, tel\deleted{s} que $\underline{r}(B)=\underline{r}(A)B$ et \deleted{que l'extens} soit $K/k$ l'extension
résiduelle. Conditions équivalentes: 1) $B$ est complet et formellement
simple sur $A$. 2) $B$ est \unsure{complet}, plat sur $A$, et $K/k$ est séparable.
3) $B$ est une algèbre de Cohen sur $A$. De plus, une extension séparable
$K$ de $k$ étant donnée, il existe une algèbre de Cohen $B$ sur $A$ se réduisant
suivant $K$, et $B$ est unique à isomorphisme de $A$-algèbres près.

La rédaction commence à faire un peu touffu, j'espère que tu pourras
la rendre plus lisible. D'ailleurs, je me répens un peu de la définition
\inserted{4.}7. telle quelle, il n'est pas vraiment naturel de supposer $B$ séparée
et complète pour sa topologie propre, mais simplement pour sa topologie
$IB$-adique (ce qui est moins fort); mais une telle modification demanderait
des réajustement analogues pour \inserted{4.}2., \inserted{4.}6. et la suite. Je t'en laisse
donc le choix, pensant que le choix adopté n'aura pas de grande importance
pratique pour la suite : dans \inserted{4.}9. à \inserted{4.}12. cela obligera simplement à
supprimer à certains endroits le mot ``complet'',
\inserted{ou \unsure{è|à} le remplacer par un autre \unsure{comple}}
\deleted{\deleted{en particulier ``Cohen'' deviendra} synonyme dans le contexte des anneaux locaux noethériens de
\deleted{``formellement} simple'', -- sous réserve de l'embêtante hypothèse de
\deleted{multiplicité} radicielle. \deleted{Il est vrai} Ce serait plutôt un avantage, \deleted{mais}
\deleted{il est vrai par contre que cela ne s'harmoniserait plus} de plus les
\deleted{anneaux de Cohen} sont}

\deleted{\uline{Corollaire \unsure{4.13}} $A\unsure{\longrightarrow|\text{---}} B$ hom loc. d'anneaux locaux
noeth, munis de leurs topologies naturelles. On
suppose $k(B)/k(A)$ à mult. \unsure{insép.} finie. Alors
Condit. \unsure{équivalentes}}

\mspage{16}
